\section{Introduction}
\label{sec:introduction}

The paradigm of the Internet of Things (IoT) has expanded into vehicular domains, establishing the Internet of Vehicles (IoV) as a critical cyber-physical framework for modern Intelligent Transportation Systems (ITS) \cite{szalay2023survey}.
Vehicular communications encompass multiple cooperative modalities, commonly categorized as Vehicle-to-Everything (V2X).
These modalities include Vehicle-to-Vehicle (V2V), Vehicle-to-Infrastructure (V2I), Vehicle-to-Pedestrian (V2P), and Vehicle-to-Network (V2N) interactions.
While traditional vehicular sensing relies strictly on autonomous on-board sensors such as radar, LiDAR, and vision systems, single-vehicle perception remains inherently limited by physical line-of-sight obstructions, environmental attenuation, and finite sensory range.
By interconnecting mobile nodes, roadside units (RSUs), and edge/cloud infrastructure, IoV enables collaborative awareness, cooperative maneuver execution, and traffic efficiency optimizations \cite{gyawali2020challenges}.

